% 相机参数笔记
% license Usr
% type Note

\subsection{画幅}
\begin{itemize}
\item 全画幅是 36×24（单位 mm，下同）
\item 索尼半画幅（APS-C）： 传感器尺寸 ≈ 23.6×15.7 （面积为全画幅的 42.8\%）
\item 佳能： 半画幅 约 22.3×14.9（38.4\%）
\item M4/3画幅： 传感器尺寸 = 17.3×13（26\%）
\end{itemize}

某画幅的\textbf{转换系数}：$\text{转换系数} = \text{全画幅宽度}/\text{该画幅的宽度}$。

\subsection{焦距}
手机一倍相机视角一般是 60° 左右。 例如 Samsung s ultra 24 的水平视角为 73.7°（$2\tan^{-1}(3/4)$）。

镜头的放大能力通常用焦距来表示（使用小孔模型计算）
\begin{equation}
\text{视角} = 2\arctan{\frac{\text{画幅长度/2}}{\text{焦距}}}~.
\end{equation}
视角越小，放大越大。 从画面长度来说，焦距和放大倍数成正比。

视野可以理解为 $2\tan(\text{视角}/2) = \text{画幅长度}/{\text{焦距}}$。

“视野为 50mm” 是一种简略的说法，意思是假设使用全画幅相机，焦距 50mm 镜头得到的视野。

\textbf{等效焦距}：意思是假设使用全画幅相机，多少焦距的镜头能达到同样的视野。 例如笔者的三星手机一倍相机就是 24mm 等效焦距。


\subsection{光圈}
光圈和焦距共同决定虚化程度。

$\text{F 值} = \text{镜头的焦距}/\text{光圈孔洞的有效直径}$

F 值越小，进光量（曝光）越大（与底片单位面积单位时间的光能成正比），同时虚化程度越高。
